
\documentclass[referee,lineno,pdflatex,sn-vancouver-num]{sn-jnl}

\usepackage{graphicx}
\usepackage{amsmath,amssymb,amsfonts,amsthm}
\input{glyphtounicode}\pdfgentounicode=1%
\usepackage{xcolor}
\usepackage{textcomp}
\usepackage{array}

\makeatletter
\let\orcidlogo\@undefined
\makeatother
\usepackage{orcidlink}

\graphicspath{{Fig/}}

\theoremstyle{thmstyleone}%
\newtheorem{theorem}{Theorem}%
\newtheorem{proposition}[theorem]{Proposition}%
\theoremstyle{thmstyletwo}%
\newtheorem{example}{Example}%
\newtheorem{remark}{Remark}%
\theoremstyle{thmstylethree}%
\newtheorem{definition}{Definition}

\begin{document}

\title[Microbiome-CDM: a provenance-traceable data model]{Microbiome-CDM: a provenance-traceable common data model and multi-program extraction method for the human microbiome literature}

\author*[1]{\fnm{Jasmeer Singh} \sur{Kalra}\orcidlink{0000-0001-7836-6860}}\email{jkalra@uni-bonn.de}
\author[2,1]{\fnm{Martin} \sur{Hofmann-Apitius}\orcidlink{0000-0001-9012-6720}}
\author[3,2]{\fnm{Marie-Christine} \sur{Simon}\orcidlink{0000-0001-6625-1265}}

\affil*[1]{\orgdiv{Bonn-Aachen International Center for Information Technology (b-it)}, \orgname{University of Bonn}, \orgaddress{\street{Friedrich-Hirzebruch-Allee 6}, \city{Bonn}, \postcode{53115}, \state{North Rhine-Westphalia}, \country{Germany}}}
\affil[2]{\orgdiv{Bioinformatics Department}, \orgname{Fraunhofer Institute for Algorithms and Scientific Computing (SCAI)}, \orgaddress{\street{Schloss Birlinghoven}, \city{Sankt Augustin}, \postcode{53757}, \state{North Rhine-Westphalia}, \country{Germany}}}
\affil[3]{\orgdiv{Faculty of Agriculture, Institute of Nutrition and Food Sciences, Nutrition and Microbiome}, \orgname{University of Bonn}, \orgaddress{\city{Bonn}, \postcode{53115}, \state{North Rhine-Westphalia}, \country{Germany}}}

\abstract{\textbf{Background:} The attributes that make microbiome sequencing
data reusable (body site, disease status, intervention, assay protocol) are
recorded inconsistently in structured submissions and survive mostly in
publication text. We present Microbiome-CDM, a method
that builds a provenance-traceable common data model (CDM) from the full text
of the 453 papers cited by a 168,464-sample gut microbiome compendium.
Separate extraction programs are kept as distinct mention sources;
mentions are normalised with two-model adjudication; findings are
catalogued with verbatim evidence; and relation triples receive a specific
label only when their sentence states the relation.

\textbf{Results:} Benchmark F1 of published recognisers predicts corpus grounding for
diseases (Spearman 0.749) but not organisms (0.023). Each of 2,505,684 mentions
and 25,788 triples is traced to its program, offsets and sentence. A blinded review found 61.8\%
(95\% CI 57.4--65.9) of sampled accepted mappings exact and 6.2\% (4.4--8.7)
wrong. After re-typing, 13,969 triples (54.2\%) carry a specific,
cue-verified relation; instead of forcing labels, the method flags 8,068
triples it cannot verify and keeps 1,210 associations
without a stated sign as unsigned. A further 4,966 specific triples derive
from catalogue findings. Against expert-curated BugSigDB, reported
directions agree in 77.0\% of determinable pairs versus a 52.0\% chance
baseline. Applied zero-shot to MiDRED's annotated microbe--disease relations,
typing reached a micro-F1 of 0.57, below fine-tuned models, but
signed answers agreed with the annotators in 88.5\%; extracted countries
agreed with GutBiomeDB's curation for the same studies in 95.8\%. Queried by genus and condition, it maps
reported directions for 22 genera in 10 conditions, which agree with
independent BugSigDB consensus in 25 of 44 comparable cells, close to chance.
On 100 sampled triples, two blinded annotators judged 70 and 79 supported or
partially supported by their sentence and agreed with each other on 87. 

\textbf{Conclusions:} A multi-program extraction method can turn the
literature behind a large microbiome compendium into a provenance-traceable
CDM. Recogniser choice has to be checked on the target corpus, because
benchmark performance transfers for diseases and not for organisms, and
abstaining where a relation cannot be verified keeps every assertion
auditable. Findings recorded only in figures and supplements are not
captured.}

\keywords{microbiome, common data model, metadata standards, FAIR, data
harmonisation, named entity recognition, large language models, text mining,
data provenance}

\maketitle

\section{Background}\label{sec:intro}

The scale of human microbiome research is now set by data aggregation rather than
data generation. The most recent large-scale synthesis integrates 168,464 samples
across 482 BioProjects~\cite{abdill2025}. The sequencing reads behind all of these
studies are downloadable, but much of what makes them analysable (host species,
body site, disease status, intervention, assay protocol) is recorded
inconsistently or not at all in structured submission metadata~\cite{goncalves2019},
and survives in usable form only in the text of the underlying publications.
Cross-study comparison at this scale is therefore bounded not by data availability
but by how much prose a project can afford to read.

Manual harmonisation of that prose is the rate-limiting step. Standardising hundreds
of papers into a shared schema takes sustained expert reading, and curated
resources such as BugSigDB show what expert curation achieves, but their scope is
fixed by curation capacity rather than by
publication rate~\cite{geistlinger2023}.

Computational resources now approach the same problem from several
directions, and each contributes something the others do not. Annotated
corpora provide ground truth: MiDRED supports knowledge-base construction for
the microbiome~\cite{hogan2024midred}, and DiMB-RE annotates 14,450 entities
and 4,206 diet--microbiome relations in 165 publications~\cite{hong2025dimbre}.
Knowledge graphs provide integration and query: MINERVA maps microbe--disease
associations across the literature with a fine-tuned language model and a
visual atlas~\cite{langarica2025minerva}; MicrobiomeKG turns the supplementary
tables of microbiome papers into scored assertions available through the
Biomedical Data Translator~\cite{goetz2025microbiomekg}; KG-Microbe harmonises
microbial traits, environments and ontologies into FAIR, AI-ready
graphs~\cite{santangelo2026kgmicrobe}; and literature mining combined with
pathway data has linked gut microbiota to drug and neurotransmitter
pathways~\cite{liu2024mikg}. Sample-level resources provide reanalysis:
GutBiomeDB curates metadata for the same 168,464 samples and adds tools for
differential and meta-analysis~\cite{zhai2026gutbiomedb}. Language models have
been used to extract and harmonise sample metadata from structured
records~\cite{ikeda2025biosample,higashi2024harmonization}, although their
reliability in biomedical text processing remains under
evaluation~\cite{chen2025llmbench}.

What these resources leave open is complementary to them: the study context
(host, body site, condition, intervention, assay, sequencing platform) and the
findings stated in the prose of the papers behind a compendium, each held with
the program that extracted it, its character offsets and its verbatim sentence.
This is the information that community reporting standards ask authors to
provide, namely the STORMS checklist for human microbiome
studies~\cite{mirzayi2021storms}, the STREAMS guidelines~\cite{kelliher2025streams}
and the MIxS/MIMARKS minimum-information checklists~\cite{yilmaz2011mixs}, and
that the FAIR principles ask to be reusable~\cite{wilkinson2016fair}. For the
large body of already published studies, it can be recovered only from text.

We built Microbiome-CDM as a pipeline rather than a single model. The pipeline
ingests the full text of each publication, converts it to section-aware structured
text, and applies an ensemble of recogniser programs (trained taggers, zero-shot
span detectors, small language models and a MeSH gazetteer) whose outputs are kept
separate instead of merged, so that every downstream row carries the identity of the
program that produced it. A guarded normalisation stage resolves mentions to
canonical terms with two-model adjudication for contested mappings, and a
schema-constrained LLM extracts entity--relation triples, each carrying a verbatim
sentence from its source paper as evidence. Because exhaustive manual annotation
of a corpus of this size is infeasible, the pipeline embeds its own evaluation:
structural checks, character-offset checks, LLM-as-judge
review~\cite{zheng2023mtbench}, a pre-specified comparison against the
expert-curated BugSigDB~\cite{geistlinger2023}, a blinded human review of a
random sample of accepted vocabulary mappings, and a blinded two-annotator check
of a random sample of relation triples.

Relative to the resources above, the contribution is the combination of four
properties in one released model. First, every mention, mapping and triple
carries the identity of the program that produced it, its character offsets and
its evidence sentence. Second, study context and findings sit in one relational
schema, so that a finding can be filtered by body site, assay or condition.
Third, the extraction abstains: a triple whose sentence does not state a
specific relation is flagged rather than given a generic one. Fourth, all 453
papers are rights-checked, so each released layer states what may be
redistributed. Table~\ref{tab:compare} sets these properties beside the
resources above.

\begin{table*}[!t]
\caption{Microbiome-CDM beside related microbiome resources. Each resource is
described by its own publication; the last column states how it and Microbiome-CDM complement each other and, where measured, how far they agree (Results).\label{tab:compare}}%
\footnotesize
\renewcommand{\arraystretch}{1.05}%
\tabcolsep=3pt%
\begin{tabular}{@{}>{\raggedright\arraybackslash}p{0.17\textwidth}>{\raggedright\arraybackslash}p{0.35\textwidth}>{\raggedright\arraybackslash}p{0.42\textwidth}@{}}
\toprule
Resource & What it provides & Complementarity with Microbiome-CDM \\
\midrule
BugSigDB~\cite{geistlinger2023} & Expert-curated differential-abundance signatures with study metadata & Reference standard for our findings (77.0\% direction agreement); we add mention-level provenance and text-only attributes \\
DiMB-RE~\cite{hong2025dimbre}, MiDRED~\cite{hogan2024midred} & Manually annotated corpora with entity and relation types; trained extraction models & Scored zero-shot: typed F1 0.25 (DiMB-RE) and 0.57 (MiDRED), below fine-tuned models; signed answers agree in 88.5\% (MiDRED). We supply corpus-scale candidates with evidence sentences \\
MINERVA~\cite{langarica2025minerva} & Literature-scale microbe--disease knowledge graph from a fine-tuned language model, with a visual atlas & Genus--condition directions agree in 72\% of comparable cells; our signed answers match 97\% of its labelled sentences. We add study-context attributes, per-row program identity and offsets \\
MicrobiomeKG~\cite{goetz2025microbiomekg} & Scored assertions from supplementary tables; available through the Biomedical Data Translator & Mines the supplements our text extraction covers least (22.1\% of supplement-sourced signatures); we cover the prose; 1 of its 42 source papers is in our corpus \\
KG-Microbe~\cite{santangelo2026kgmicrobe} & FAIR, ontology-grounded graphs of microbial traits and environments & Trait data for 86\% of our genera; its Disbiome directions agree with ours in 72\% of cells; we supply study-level context and findings \\
GutBiomeDB~\cite{zhai2026gutbiomedb} & Curated metadata and analysis tools for the 168,464 compendium samples & Same compendium, different layer: samples there, the papers here; our extracted country agrees with its curation in 95.8\% of studies, and we cover 77.1\% of its disease labels \\
PubTator3~\cite{wei2024pubtator3} & Entity and relation annotation across PubMed and PMC & One of our mention sources; we keep its output beside other programs rather than merged \\
LLM metadata extraction~\cite{ikeda2025biosample} & Language-model extraction of terms from structured sample records & We apply the same idea to full-text papers, with evidence offsets and a judge-based audit \\
\midrule
Microbiome-CDM (this work) & 453 rights-checked full-text papers; 2.5 million mentions from separate programs; mappings; catalogue; 25,788 triples, each traced to program, offsets and sentence & Strengths: provenance at row level, explicit abstention, licence discipline. Gap: small human checks (500 mappings; 100 triples, two annotators), no annotated gold standard \\
\botrule
\end{tabular}
\end{table*}

Here we describe the pipeline and the resulting resource. Build~3.0.3 of the
resource holds
162,409 canonical terms, 334,231 candidate mappings of which 314,772 pass
vocabulary guards and 18,272 are accepted by two-model agreement, and 25,788
canonical entity--relation triples, each carrying a sentence from its
source paper as evidence (for 24,157 of them located in the section text
verbatim or after normalisation). As a check on the extraction, an automated judge
screen by a single large language model (DeepSeek-V4.1-Flash, accessed through
the Fraunhofer Genie gateway) of the 16,032 triples whose cited sentence is CC BY
or CC0 licensed
classifies 56.0\% as supported, 22.6\% as partially supported, 18.2\% as not
supported and 3.3\% as conflicting; these are model judgements, reported as
concordance measures rather than accuracy claims. In the following sections, we
detail the multi-phase processing lineage, from corpus acquisition to the
common data model, and present the computational evaluation of the pipeline.

\section{Methods}\label{sec:methods}

\subsection{Overview}\label{sec:overview}
\begin{figure}[htbp]
  \centering
  \includegraphics[width=\linewidth]{fig_overview.pdf}
  \caption{Pipeline overview. Five stages build the common data model
  from the literature underlying a human gut microbiome compendium: (1)
  assembly and licence-checking of the cited corpus (453 BioProjects,
  447 publications, 4,006 sections); (2) multi-program entity extraction
  from kept-separate sources (trained and zero-shot recognisers, two
  small-LLM runs, a MeSH dictionary), yielding 2,505,684 mentions with
  character positions; (3) normalisation to 162,409 canonical terms and
  guarded ontology mapping (334,231 candidates, 314,772 kept, 18,272
  accepted by two-model agreement); (4) the CDM 3.0.3 relational build
  (catalogue records, evidence quotes, 25,788 relation triples); and (5)
  frozen, versioned releases with checksummed manifests. The evaluation
  band underneath is mostly computational (structural checks, evidence
  checks, model judges, curated-reference comparisons, NER benchmarks
  and research-question use cases), complemented by blinded human reviews
  of 500 sampled vocabulary mappings and 100 sampled relation triples.}
  \label{fig:overview}
\end{figure}

Four design choices run through all five stages
(Figure~\ref{fig:overview}).
\textbf{Provenance first:} every row produced downstream carries the
identity of the extraction program, the source paper, section and
character offsets, so any value can be traced back to its source
text. We record this provenance directly as columns rather than mapping it
to the W3C PROV-O ontology: the question an auditor asks here is which
program produced which row from which sentence, and these columns answer it
without a separate provenance graph. \textbf{Sources kept separate:} the four mention sources are
never merged; per-program identity is preserved into the CDM, which
allows later re-weighting without re-extraction.
\textbf{Mostly computational evaluation:} exhaustive manual annotation
of a corpus of this size is infeasible, so evaluation relies on
structural checks, model judges and concordance against curated
references, reported as concordance rather than accuracy, and is
anchored by blinded human reviews of random samples of accepted
vocabulary mappings and of relation triples.
\textbf{Licence-aware reporting:} every paper is licence-checked, and
results that depend on source sentences are reported only for papers
licensed CC~BY or CC0. One screen run outside our computing cluster (DeepSeek-V4.1-Flash on Genie,
a model service hosted by Fraunhofer, one of the authors' institutions)
received its sentences
before this filter was applied, including 9,629 sentences from papers
under other licences; its results are reported for the CC~BY/CC0 subset
only. Later external screens received either term and concept labels only
or, to resolve unsigned associations, text from CC~BY or CC0 papers only.

\subsection{Stage 1: Corpus and sectioning}\label{sec:stage1}

The foundational corpus was constructed to represent the literature
underlying large-scale human microbiome data aggregation. We seeded the
corpus acquisition using BioProject accessions~\cite{barrett2012bioproject}
(PRJDB, PRJNA, PRJEB) and digital object identifiers (DOIs) of the
literature cited in a recent synthesis of 168,464 samples~\cite{abdill2025}.
Automated ingestion of PDF and TXT formats was managed via a custom
processing script (\texttt{human\_microbiome\_corpus\_builder.py}), which
applied metadata heuristics and topic classification to yield an initial
corpus of 453 reference papers.

To ensure data integrity, a deduplication stage (stage 30) isolated 447
canonical papers and excluded 6 duplicates. Rather than treating full texts
as monolithic blocks, a section-aware parsing algorithm (stage 40) segmented
the canonical papers into 4,006 distinct sections, comprising 26.9 million
characters. Subsequent full-text validation independently cross-checked 296
of these papers against external bibliographic metadata. One paper
(HMPAPER\_000372) was retrieved only as a publisher redirect page (909
characters, no article text) and yields no catalogue records; it is kept in
the 453 for completeness, and analyses that need article text use the other
452.

\subsection{Stage 2: Entity extraction with multiple programs}\label{sec:stage2}

To maximise entity recall without sacrificing provenance, we designed an
extraction pipeline that treats different named entity recognition (NER)
programs as independent annotators~\cite{patel2026microbialner}. The pipeline
deployed an ensemble of recogniser programs, whose outputs were intentionally
kept separate rather than merged, ensuring every downstream row retains the
exact identity of the program that produced it.

The extraction ensemble consisted of four distinct programmatic branches:
\begin{enumerate}
\item \textbf{CPU Section Ensemble:} a baseline configuration
(\texttt{zenodo\_v1}) comprising 7 discrete models, which identified
1,705,449 entity rows in the CDM across 2,817 body sections (the terminology
release lists 1,705,501 before loading).
\item \textbf{Small Open LLM (Section-level):} Qwen3.5-4B~\cite{qwen35card} on section-level inputs, generating 186,024 entity rows.
\item \textbf{Small Open LLM (Full-body):} Qwen3.5-4B deployed via
vLLM~\cite{qwen35card,kwon2023vllm} on full-text inputs, yielding 366,123
rows.
\item \textbf{MeSH Gazetteer:} a rule-based 2026 dictionary
search~\cite{lipscomb2000mesh} contributing 248,088 rows.
\end{enumerate}

All mentions were carried forward with precise character offsets linking
them to their verbatim source sentences. Of the 2,505,684 mentions recorded
in the final model, 2,505,628 are anchored at their offsets; a systematic
offset shift in 75 sections was detected and corrected, and all but 56
mentions reproduce their source surface form at the recorded offsets.

\subsection{Stage 3: Normalisation and vocabulary mapping}\label{sec:stage3}

Because raw NER outputs contain high orthographic variation, the pipeline
executed a multi-step terminology normalisation chain (T1--T5). Following
the assembly of the 2.5 million mentions (T1) and the extraction of their
contextual definitions (T3), a two-level canonicalisation process (T2)
collapsed textual variants, yielding an initial space of 166,012 canonical
terms, which was subsequently refined to 165,991 terms.

To bridge these terms to standardised ontologies (e.g.\
UMLS~\cite{bodenreider2004umls}, OMOP~\cite{hripcsak2015ohdsi}, NCBI
Taxonomy~\cite{schoch2020taxonomy}, BioPortal~\cite{whetzel2011bioportal}),
we implemented a guarded mapping strategy (T5). The pipeline generated
334,231 candidate mappings. Hard programmatic guards filtered this set
down to 314,772 biologically plausible mappings. To resolve the most
contested conceptual mappings, the pipeline used a two-model adjudication
framework (DeepSeek-R1~\cite{guo2025deepseekr1} and
Gemma~\cite{gemma2026gemma4}). Requiring both models to agree accepted
18,272 mappings, with the agreed relation (exact, broader, narrower or related)
recorded for each. Finally, a draft taxonomy
(T4) was exported, organising the terms into 16 upper categories and 66
subcategories, forming Terminology release 1.0.1.

For build 3.1.0 the vocabulary was curated further without altering the
3.0.3 tables. Concept identifiers that spell the same code differently
(for example the two MeSH IRI forms) were first collapsed
deterministically, reducing 13,330 accepted identifiers to 12,657 concepts.
Terms and concepts were then deduplicated in passes of increasing
tolerance: exact matching after normalisation, fuzzy matching restricted
to harmless single-word differences, expansion of abbreviated genus names
when exactly one genus in the corpus fits, and an embedding pass that
merges two items only if both their names (SapBERT~\cite{liu2021sapbert}) and their contextual
definitions (MiniLM sentence embeddings~\cite{liu2019sbert,minilmcard}) agree and the wording difference is explained by word
order, plural forms or an acronym. Guards block any merge across entity
types or between different numbers, single letters, signs, opposite
prefixes or accepted concepts, and two organism names are merged only when
they denote the same taxon, so that, for example, \textit{Lactobacillus}
and \textit{Lactococcus} remain distinct. Deduplication reduced the
162,409 canonical terms to 157,815 and the concepts to 12,182. A relevance
filter (DeepSeek-V4.1-Flash, given only each label and its entity type)
then removed strings that are not concepts, such as author names, figure
references and sentence fragments, or that lie outside biomedical
research, keeping 116,755 terms, 11,912 concepts and 15,769 accepted
mappings between them.

Two further layers type what extraction left generic. First, each of the
53,780 kept terms that carried only the generic ENTITY type was re-typed by
a locally run language model (Qwen3.5-9B~\cite{qwen35card}, 4-bit AWQ
weights~\cite{lin2023awq}, served with vLLM on a single GPU), given the
term, its merged variants, its model definition and up to three corpus
sentences from distinct papers. A JSON schema restricts the answer to the
14 CDM types, to four proposed extension types for things the 14 cannot
hold (anatomy or body site, population or host group, measure or unit, and
geography), or to NOT\_A\_CONCEPT for words that name no scientific thing.
Second, each of the 14,867 triples that the repair pass marked as stating
no specific relation was re-typed (of the 17,481 triples without a specific,
cue-verified relation after repair, the remaining 2,614 were already flagged
as unsupported or carried a specific label without a verified cue) from up to three of its source sentences. The
model first quotes the words that relate subject and object and only then
chooses a relation, from the existing closed list extended by signed
association labels (positively or negatively associated, differentially
abundant, biomarker of, predicts, increases or decreases over time) or one
of three exits: association without direction, null finding, or not
supported. A relation counts as specific only when its quoted cue is found
verbatim in a source sentence; a null finding also requires negation
wording in the cue, and a reversed direction is overruled when the
sentence reads subject, cue, object in active voice. Both prompts were
frozen after constructed test sets on which 19 of 20 terms and 17 of 17
sentences were typed as expected; the prompt for catalogue-derived triples
(below) typed 18 of 20 constructed cases as expected.

Two further passes completed the relation layer. The 1,708 triples that
still stated only an unsigned association were read again with more
context: first the evidence sentence and its neighbours and, if no sign or
direction was found, the full section that holds the sentence. Text from
CC~BY or CC0 papers was sent to DeepSeek-V4.1-Flash on Fraunhofer's Genie
service (959 triples); all other triples were processed by the local
Qwen3.5-9B (749 triples), so that their text did not leave our computing
cluster. A signed label was kept only if its verbatim cue contained sign
wording other than ``risk'' and, for the section step, named at least one
of the two entities, and each answer was checked by a separate yes/no
verification. Finally, 13,046 catalogue findings whose evidence quote is
anchored in the paper and that name a taxon or feature, a comparison and a
direction were converted into candidate triples when the quote contains
the subject, the object and wording for the direction (7,424 candidates,
6,856 of them not already in the model). The local model typed each
candidate from its quote in section context, under the same checks and
verification.

\subsection{Stage 4: Data catalogue and common data model}\label{sec:stage4}

To synthesize isolated entities into contextual biological knowledge, we
executed a schema-constrained relation extraction, following comparable
LLM-based extraction pipelines in microbiology and protein
research~\cite{lu2025microllm,chen2026buffers}. A secondary LLM pass
(CATNER v2 utilising Qwen3.5-4B via vLLM) scanned the processed text to
extract complex findings, assays, cohorts, samples, and datasets. This was
followed by a deterministic cleanup (Phase A) and an LLM-assisted linking
phase (Phase B). To address missing vocabulary and unmatched relations, the
pipeline triggered gap-repair chains (G1--G10), which recorded 40,141
per-term type adjudications reconciling the entity types assigned by the
different recognisers with the canonical terminology, and linked catalogue
relations directly to the canonical terms.

The finalised architecture (Microbiome-CDM 3.0.3) was compiled into a highly
structured SQLite relational database (Figure~\ref{fig:cdm}). The migration to the final build
(3.0.3-rc1 to frozen) was strictly additive, generating a reproducible schema
of 44 tables containing 8,490,496 rows. The final model successfully
encapsulated 2,505,684 mentions, 162,409 canonical terms with their guarded
vocabulary mappings, and 25,788 canonical entity--relation triples, forming
a graph of 22,579 unique nodes.
\begin{figure}[htbp]
  \centering
  \includegraphics[width=\linewidth]{fig3_cdm_architecture.pdf}
  \caption{Architecture of CDM 3.0.3. The model resolves 33 entity
  types across 10 modules covering source documents, document
  structure, terminology, semantic mappings, processing provenance,
  quality review, evidence, findings, research context and catalogue
  releases. Per-entity row counts and population status are shown
  (fully populated, partially populated, or embedded in another
  table). Evidence records anchor to their document-structure spans,
  and finding and triple records link to canonical terms; every
  record additionally carries its source paper and extraction run,
  which together with the checksummed releases allow any derived
  value to be traced back to source text.}
  \label{fig:cdm}
\end{figure}

\subsection{Stage 5: Computational evaluation and human spot check}\label{sec:stage5}

Because the 453-paper corpus size precludes comprehensive human
ground-truth annotation, the pipeline embeds a mostly computational
evaluation framework, complemented by two small blinded human checks. We
evaluated the pipeline along the following dimensions:

\begin{enumerate}
\item \textbf{Structural Integrity:} automated scripts verified
cryptographic hashes, text offsets and identifier linkage across the 44
SQLite tables, with 118 of 121 tests passing (2 failures, 1 conditional).

\item \textbf{Metadata Concordance:} we conducted a source-first crosscheck
using deterministic matching, successfully recovering 88.7\% of accession
numbers, 89.2\% of sequencing platforms and 72.1\% of 16S rRNA gene identifiers
natively present in the text.

\item \textbf{External Curation Comparison:} we benchmarked our automated
triple extraction against expert-curated databases, specifically testing
against BugSigDB~\cite{geistlinger2023} (using 47 overlapping papers),
Disbiome~\cite{janssens2018disbiome}, and CTD~\cite{davis2026ctd}.

\item \textbf{NER Benchmarking:} the underlying NER programs were evaluated
against public gold-standard datasets (BC5CDR and BioRED) using 27 distinct
configurations. PubTator3 achieved an exact F1 of 90.1 on BC5CDR and
HunFlair2 achieved 80.6 on BioRED, with the best small LLM (Qwen3.5-4B,
3-shot) reaching 69.5 and 66.2 respectively~\cite{wei2024pubtator3,
sanger2024hunflair2,qwen35card}.

\item \textbf{Evidence-Anchored Screen:} to assess the logical validity of
the 25,788 extracted triples, we utilised an
LLM-as-judge~\cite{zheng2023mtbench} (DeepSeek-V4.1-Flash, accessed via the
Fraunhofer Genie gateway). The model evaluated each triple against its own
verbatim anchor sentence, assessing whether the predicate is justified by
the sentence, whether the polarity is correct, and whether the quoted
evidence matches the text, yielding an overall support verdict. To adhere
strictly to licensing constraints for external API processing, we report
results only for the 16,032 triples sourced from CC BY or CC0 licensed
papers. The screen classified 56.0\% of these relations as fully supported
by the text, 22.6\% as partially supported, 18.2\% as unsupported, and
3.3\% as conflicting. We report these strictly as model concordance
judgements, not as absolute accuracy
claims~\cite{chen2025llmbench,deepseekv41card}.

\item \textbf{Catalogue Judges:} three locally run model judges
(Qwen3.5-9B, Llama-3.1-8B and Mistral-7B~\cite{qwen35card,grattafiori2024llama3,jiang2023mistral})
checked catalogue values under a decomposed rubric. Links between records
were checked structurally rather than by a model. For classification
fields each judge answered the question blind, and its answer was compared
with the stored value; for stated values a judge had to quote an evidence
span, which was accepted only if found in the source passage and covering
the value; record-level summary labels were rated as fair or unfair.
Verdicts were constrained by a JSON schema. The rubric was developed on a
1\% development split (536 records, assigned by hash of the record
identifier) and frozen before the 51,837 held-out records were judged.
Agreement is reported as Cohen's and Fleiss' $\kappa$~\cite{cohen1960kappa},
Gwet's AC1 and per-verdict specific agreement, as agreement between models
rather than accuracy.

\item \textbf{Two-annotator triple check:} two annotators with microbiome
expertise independently judged a simple random sample of 100 of the 16,032
screened triples (CC BY or CC0 sentences only), using the same SUPPORTED,
PARTIAL, NOT\_SUPPORTED and CONFLICTS definitions as the model judge and
seeing only the triple and its sentence. The sample, protocol, blank sheets
and model verdicts were fixed and checksummed before annotation, and model
verdicts were opened only after both sheets were complete. Agreement is
reported as observed agreement and Cohen's $\kappa$ for the four labels and
for supported or partially supported versus not.

\item \textbf{Gold standards and curated resources:} the relation-typing and
entity prompts were applied zero-shot, unchanged and locally (Qwen3.5-9B), to
the human-annotated DiMB-RE~\cite{hong2025dimbre} and MiDRED~\cite{hogan2024midred}
corpora and to 1,100 manually labelled microbe--disease
sentences~\cite{karkera2023microbiomere} distributed with
MINERVA~\cite{langarica2025minerva}, with gold entity pairs given and our labels
mapped to each corpus's label set. Study metadata were compared per BioProject
with GutBiomeDB's curated metadata~\cite{zhai2026gutbiomedb}, and
genus-by-condition directions with MINERVA's graph and with the Disbiome
associations~\cite{janssens2018disbiome} contained in
KG-Microbe~\cite{santangelo2026kgmicrobe}; paper and genus overlap were measured
for MicrobiomeKG~\cite{goetz2025microbiomekg} and KG-Microbe. Competitor data
were used under their own terms and processed only locally.
\end{enumerate}

\paragraph{Statistical reporting.} Every percentage is stated with its
denominator or the count it is a share of, and shares are computed over
distinct items unless stated otherwise (mention-weighted validity is labelled
as such). Proportions from sampled reviews and judge agreement carry 95\%
Wilson score intervals~\cite{wilson1927}. Agreement with curated references is
tested against a chance level, either the same comparison made across
different papers or the base rate of the compared values, by an exact
one-sided binomial test. Kept and dropped triples are compared by Fisher's
exact test with an odds ratio and a Newcombe interval for the difference, and,
because triples are nested in papers, also by a paper-cluster bootstrap;
differences in benchmark F1 are tested by a paired document bootstrap with
5,000 resamples~\cite{efron1993bootstrap}. $p$-values are Holm-adjusted within
each family of related comparisons, and results that do not survive
adjustment are described as trends. The analyses are descriptive; no
sample-size calculation applies, because the corpus is fixed by the
compendium.

\section{Results}\label{sec:results}
\begin{table*}[!tbp]
\caption{Summary of the extraction pipeline and of the two frozen builds.
Quantities are given per stage and are reproducible from the deposited
release; \texttt{CL} rows are the CDM 3.1.0 layers and \texttt{JV} rows the
re-typing layers and the catalogue judges.
Percentages and judge-derived values are concordance or model-judgement
measures, not accuracy against human annotation.\label{tab:summary}}%
%% GENERATED by src/scripts/figures/make_pipeline_summary_table.py from the claims register -- do not edit;
%% re-run the script. Row labels: author review. Caption: written by the authors (in manuscript.tex).
\footnotesize
\renewcommand{\arraystretch}{0.95}%
\tabcolsep=3pt%
\begin{tabular}{@{}l>{\raggedright\arraybackslash}p{0.40\textwidth}>{\raggedright\arraybackslash}p{0.27\textwidth}l@{}}
\toprule
Stage & Quantity & Value & Register \\
\midrule
Corpus & papers / distinct publications & 453 / 447 & C01, C02 \\
 & sections (SHA-256 hashed) & 4,006 & C07, C08 \\
 & compendium samples in corpus BioProjects & 156,330 of 168,464 (92.8\%) & C05, C06, C09 \\
 & papers whose text may be redistributed & 274 of 453 & RT01, RT02 \\
\midrule
Extraction & mentions, four sources kept separate & 2,505,684 & M01 \\
 & CPU NER ensemble / Qwen3.5-4B full-body / Qwen3.5-4B section-level / MeSH & 1,705,449 / 366,123 / 186,024 / 248,088 & M03--M06 \\
\midrule
Normalisation & canonical terms (after review) & 162,409 & T02 \\
 & vocabulary mappings: candidates / kept by guards / accepted by two-model agreement & 334,231 / 314,772 / 18,272 & P01, P02, P05 \\
\midrule
Catalogue & evidence-quoted records & 52,400 & K01 \\
 & quotes found verbatim at their offsets & 88.2\% & K08 \\
\midrule
CDM 3.0.3 & tables / columns & 44 / 866 & X01, X02 (CDM) \\
 & registered extraction runs & 45 & X03 (CDM) \\
 & relation triples / with a generic relation label & 25,788 / 17,717 & R01, R02 \\
 & schema, value and provenance tests: pass / fail / conditional & 118 / 2 / 1 & X07 \\
\midrule
CDM 3.1.0 & terms after dedup / kept after relevance filter & 157,815 / 116,755 & CL01 \\
 & accepted concepts after ID normalisation / dedup / kept & 12,657 / 12,182 / 11,912 & CL02 \\
 & accepted mappings, clean rows (of which EXACT agreements in 3.0.3) & 15,769 (13,540) & CL03, CL04 \\
 & triples with a specific relation and verified cue & 8,307 & CL05 \\
 & triple tiers 1 / 2 / 3 & 1,794 / 6,263 / 17,731 & CL06 \\
 & terms with a model definition: contextual / general / descriptive & 120,091 / 530 / 41,779 & CL08 \\
\midrule
Evaluation & best exact F1, public gold standards (BC5CDR / BioRED) & 90.1 / 80.6 & B01, B02 \\
 & BugSigDB: same direction of reported findings, of n pairs (chance) & 77.0\% of 261 (52.0\%) & N05 \\
 & triples judged against a CC BY / CC0 sentence (LLM screen) & 16,032 & W02 \\
 & judged supported / partial / not supported / conflicting & 56.0\% / 22.6\% / 18.2\% / 3.3\% & W03 \\
\botrule
\end{tabular}

\begin{tablenotes}%
\item LLM-screen and BugSigDB values are concordance or model judgement, not accuracy.
\end{tablenotes}
\end{table*}

\begin{figure}[htbp]
  \centering
  \includegraphics[width=\linewidth]{fig4_results.pdf}
  \caption{Computational evaluation.
  (a)~Share of distinct entity strings judged valid and correctly
  typed, per extraction configuration (string screen, distinct
  strings). (b)~Agreement between catalogue metadata and the
  expert-curated BugSigDB over the 47 shared papers, against
  cross-paper chance baselines; daggered fields are matched by
  SapBERT semantic similarity rather than exactly.
  (c)~Catalogue coverage of BugSigDB signatures by where BugSigDB
  recorded the signature. (d)~BugSigDB support for taxon triples
  kept versus dropped by model-based cleaning. All panels report
  concordance or judged validity, not accuracy against a human gold
  standard.}
  \label{fig:results}
\end{figure}

The five stages produced CDM build 3.0.3, a relational model of 44
tables and 866 columns in which every extracted value carries paper,
section and character-offset provenance. Of 121 schema, value and
provenance tests, 118 pass, two fail and one is conditional, and no row
names an unregistered extraction run. Seventeen entity types are
fully populated, fifteen partially (five of them embedded in other
tables' rows), and one is planned. The build, terminology and
catalogue releases were frozen with SHA-256 manifests
(Section~\ref{sec:stage4}). Build 3.1.0 leaves the entity tables unchanged and
adds the curated layers described in Section~\ref{sec:stage3}: 116,755 kept
terms, 11,912 kept concepts and 15,769 accepted mappings between them, of
which 13,540 of the original accepted rows are exact agreements and the
remainder broader, narrower or related concepts. After re-typing, no kept
term has the generic type: 103,088 (88.3\%) carry one of the 14 CDM types,
8,856 a proposed extension type and 4,811 are flagged as not a concept.
A blinded reviewer
classified 307 of 497 randomly sampled accepted mappings (61.8\%, 95\%
Wilson interval~\cite{wilson1927} 57.4--65.9\%) as exact, 159 as broader, narrower or related
(32.0\%, 28.0--36.2\%) and 31 (6.2\%, 4.4--8.7\%) as wrong, judged 485 of
496 terms (97.8\%, 96.1--98.8\%) to be genuine concepts and 426 of 493
entity types (86.4\%, 83.1--89.2\%) to be correct; the second reviewer's
100-mapping overlap is pending. Overall, 466 of 497 mappings (93.8\%,
91.3--95.6) were not wrong. Exactness depended strongly on how much support a
term had: it rose from 39.7\% (96 of 242) for terms mentioned once or twice to
89.6\% (60 of 67) for terms mentioned more than 20 times, and from 31.9\% (45 of
141) for terms found by one extraction program to 88.4\% (38 of 43) for terms
found by six; by vocabulary it ranged from 41.7\% (NCIt) to 91.2\% (ChEBI).
Because mention counts and program counts are released with every term, users
can restrict the mappings to well-supported terms. Table~\ref{tab:summary} sets
out the five stages and the two frozen builds, with the source of every value.

Deterministic matching locates 70.5\% of catalogue text values and
75.4\% of numeric counts inside their own quote-anchored evidence
passage; the remainder are dominated by coded
values that cannot be string-checked.
Three local model judges then checked 51,837 held-out catalogue records,
and all 167,197 links between records point to existing records of the
same paper. The two most consistent judges (Llama-3.1-8B and Mistral-7B)
agree on 82.9\% of 442,772 field verdicts (95\% Wilson interval
82.8--83.0; Cohen's $\kappa$ 0.566, Gwet's AC1 0.788), and when one fails a
value the other fails it too in most cases (specific agreement 0.598). With
sampling error this small the informative quantity is the variation across
fields, not the interval. Fleiss' $\kappa$ over all three judges is
0.471, 0.535 for stated values and 0.616 for the judges' blind answers to
classification questions. The remaining disagreement sits in a few fields
where the models default to different answers, for example whether a
control-group record describes one group or a case--control design. We
therefore use the judges as screening tools rather than raters: a majority
fails 13.1\% of stated values (33,253 of 254,022; 95\% Wilson interval
12.96--13.22\%), concentrated in DNA
extraction method (68.9\%), sequencing platform (56.3\%), body site
(40.4\%) and intervention agent (40.2\%). A pattern-based extractor containing no model component finds
its accession numbers, sequencing platforms and 16S rRNA gene regions present
in our catalogue in 88.7\%, 89.2\% and 72.1\% of cases, an agreement
measure between two methods, not a recall bound.

In the large string screen (730{,}195 judged items), distinct-string
validity separates the program families clearly
(Figure~\ref{fig:results}a). The MeSH gazetteer reaches 67.3\% of distinct strings
judged valid and correctly typed, BioBERT-diseases~\cite{lee2019biobert} 60.9\% and the
trained BERT chemical tagger 59.3\%; the best zero-shot recogniser,
GLiNER v2.5 medium, reaches 43.6\%~\cite{zaratiana2024gliner,zaratiana2025gliner2}; the
small language models span 15.7\% (Llama zero-shot) to 35.5\% (Gemma
zero-shot), with Qwen3.5-4B at 34.5\% (full-body run) and 36.6\% (section-level
run), or 50.7\% when mentions are weighted instead of distinct
strings. Validity also depends on CDM type: chemicals (44.3\%), data
analysis (43.1\%) and diseases (42.8\%) screen far better than
experimental methods (14.1\%) or study design (16.1\%). For the
build-3.0.3 triples, the same screen judges 39.3\% to be well-formed
claims, with valid subject and object in 38.7\%. After relation repair and
re-typing, 13,969 of the 25,788 triples (54.2\%) carry a specific
relation whose cue phrase is found in the source sentence; 8,068 are
flagged because the pipeline could not verify them against their sentence,
1,604 carry a specific label
whose cue was not found verbatim, 1,210 state an association without
direction even when the whole section is read, 924 report a null finding
and 13 have a vague endpoint. Reading more of the paper resolved 405 of
the 1,708 unsigned associations, 85 of them only from the wider section; 1,210
stayed unsigned, and the other 93 moved to another exit or received a specific
label whose cue was not found.
The catalogue's findings add 4,966 specific, cue-verified triples from 430
papers and 848 null findings, so that 18,935 of 32,644 triples (58.0\%)
carry a specific, verified relation. We treat the remainder as abstention,
not as relations delivered: where the sentence does not state a specific
relation, the triple is flagged or kept unsigned rather than given a generic
label, so that every typed relation can be checked against its cue phrase.

In the blinded two-annotator check of 100 triples, judged as originally
extracted, the annotators agreed on
whether a triple was supported or partially supported by its sentence for 87
(87.0\%, 95\% Wilson interval 79.0--92.2; Cohen's $\kappa$ 0.66) and on all four
labels for 64 ($\kappa$ 0.46); most disagreements separated SUPPORTED from
PARTIAL for hedged sentences. Annotator A judged 35 triples supported, 35
partially supported and 30 not supported; annotator B 65, 14 and 20 (one
conflicting). The DeepSeek-V4.1-Flash screen agreed with the annotators about
as well as they agreed with each other ($\kappa$ 0.34 and 0.42 on four
labels). In an exploratory comparison with the final relation layer, 33 of 42
triples that the pipeline types as specific and cue-verified were judged
supported or partially supported by annotator A (78.6\%, 64.1--88.3) and 38
by annotator B (90.5\%, 77.9--96.2). Of the 41 triples the pipeline flagged
as unverified, annotator A judged 25 and annotator B 30 to be supported or
partially supported: the flag is conservative, and many flagged triples state
a real relation that the relation vocabulary does not cover. Verified and
flagged triples therefore differ in support less than their labels suggest
(78.6--90.5\% versus 61.0--73.2\% supported or partially supported); the cue
check confirms that a specific label is worded in the sentence, not that the
remaining triples are wrong. Expanding the closed relation vocabulary is the
most direct way to recover flagged relations in a future build. Agreed labels
from the annotators' adjudication are pending.

On public gold standards~\cite{li2016bc5cdr,luo2022biored}, spanning
27 configurations of 21 systems, the best systems are PubTator3 with
exact F1 90.1 on BC5CDR~\cite{wei2024pubtator3} and HunFlair2 with
80.6 on BioRED~\cite{sanger2024hunflair2}; the best small language
model (Qwen3.5-4B, three-shot prompting) reaches 69.5 and
66.2~\cite{qwen35card}. Per BioRED type, the best disease, gene/protein
and organism systems are {\em not} trained on BioRED (HunFlair2:
86.6, 83.3 and 94.4), while the best chemical, variant and cell line
systems are trained in-domain (PubTator3: 88.8, 82.3 and 80.8).
These rankings transfer to the microbiome corpus only partially:
Spearman correlation between BioRED F1 and corpus grounding is 0.749
for diseases but 0.023 for organisms (0.2 for genes and proteins, 0.252
for chemicals, 0.112 for variants), and no configuration of the 26
corpus benchmark configurations grounds more than 41.3\% of its own
extracted strings. Public benchmark F1 therefore cannot be read as
corpus-level accuracy, and program choice has to be re-checked against
the target corpus rather than inherited from benchmark leaderboards.

Against the expert-curated BugSigDB~\cite{geistlinger2023}, over 47
shared papers and 270 signatures, 27.7\% of BugSigDB taxa are covered
by our findings (exact or same genus) against a 9.0\% cross-paper
chance level, and 23.5\% by exact NCBI taxon ID; conversely, 57.3\% of
our taxa also appear in BugSigDB signatures against 18.1\% chance.
Where a direction comparison is determinable, reported directions
agree in 77.0\% of 261 taxon pairs, against a 52.0\% base rate.
Coverage tracks where BugSigDB recorded the signature: every
text-sourced signature is covered, versus 51.4\% of table-sourced,
27.6\% of figure-sourced and 22.1\% of supplement-sourced ones
(Figure~\ref{fig:results}c). Study metadata agrees closely for 16S rRNA gene
region (97.7\%), country (94.0\%), condition (92.7\%, compared by
SapBERT), body site (84.2\%) and statistical test (85.7\%), and less
well for group size (56.2\%) and group name (54.2\%).

\paragraph{Reuse example: reported directions by genus and condition.}
As a first reuse of the model, we asked a question a microbiome researcher
would ask of it: which genera do the papers report as increased or decreased
in each condition? Querying the catalogue findings by NCBI genus and by
condition returns, for the 22 most frequently reported genera and the 10 most
frequent conditions, the number of papers reporting each direction, each
traceable to its sentence (Figure~\ref{fig:taxon}). The map recovers
well-known patterns, for example \textit{Faecalibacterium} reported as
decreased in inflammatory bowel disease (IBD) and colorectal cancer, and
\textit{Fusobacterium} as increased in colorectal cancer. To compare it with
an independent source, we took BugSigDB signatures from papers outside our
corpus and kept the 44 cells with a clear consensus (at least three
signatures, at least two thirds in one direction). Our corpus reports the
same direction in 25 of them (56.8\%; 95\% Wilson interval
42.2--70.3). This is close to chance and is not a
validation; the 19 disagreements mark genus--condition pairs on which the
literature itself is divided, and the map is a reported-findings summary,
not a meta-analysis.

\begin{figure}[htbp]
  \centering
  \includegraphics[width=0.82\linewidth]{figF_taxon_condition.pdf}
  \caption{Reported direction per genus and condition, queried from the
  catalogue findings. Colour gives the balance of papers reporting the
  genus as increased (red) or decreased (blue) for the condition; the number
  is the count of papers. A check mark or cross shows agreement or
  disagreement with the independent BugSigDB consensus (signatures from
  papers outside our corpus; at least three signatures, at least two thirds
  in one direction); 25 of 44 comparable cells agree. Reported findings,
  not a meta-analysis.}
  \label{fig:taxon}
\end{figure}

\paragraph{Gold standards and curated resources.} Applied zero-shot to the
human-annotated DiMB-RE corpus with gold entity pairs given, our relation
typing assigned the exact label and direction to 37.7\% of 2,410 gold
relations (95\% Wilson interval 35.8--39.7), and its typed extraction F1 was
0.25 (document bootstrap 0.23--0.27), below the published fine-tuned
end-to-end model (0.445); most errors were relations asserted for entity pairs
the annotators left unrelated (32\% of 10,203). Entity recognition reached an
F1 of 0.69 (0.67--0.71, overlapping spans) on DiMB-RE types with a CDM
equivalent, with 90.9\% recall for microorganisms. On MiDRED's 3,116 doubly
annotated microbe--disease relations, four-class micro-F1 was 0.57
(0.55--0.59; test split 0.55, against 0.72 for GPT-4 in context and 0.91 for a
fine-tuned model); when both the gold label and our answer were signed, they
agreed in 1,920 of 2,169 cases (88.5\%, 87.1--89.8), and our answer was ``no
relation'' for 23.2\% of true relations. On 1,100 manually labelled
microbe--disease sentences, our signed answers agreed with 449 of 463 signed
labels (97.0\%, 95.0--98.2), with abstention on 47.1\% of signed sentences.
The zero-shot steps are therefore precise when they commit to a sign but
cautious, and below fine-tuned extractors in recall and overall F1.

Against curated metadata, our text-extracted fields agreed with GutBiomeDB for
the same BioProjects on country in 205 of 214 studies (95.8\%, 92.2--97.8) and
covered 229 of 297 curated disease labels (77.1\%, 72.0--81.5); disease
matching was done by the same local model, so this is a concordance measure.
Genus-by-condition directions agreed with MINERVA's graph in 95 of 132
comparable cells (72.0\%, 63.8--78.9; chance about 50\%), with most
disagreements on probiotic genera that intervention studies report as
protective, and with the curated Disbiome associations in KG-Microbe in 53 of
74 (71.6\%, 60.5--80.6); both sources may share papers with our corpus, so
these comparisons are not independent. MicrobiomeKG shares one of its 42
source publications with our corpus and contains 26.8\% of our 321 catalogue
genera, whereas KG-Microbe holds trait data for 86.0\% of them, so these graphs
complement rather than duplicate the CDM.

For relation triples, a licence-filtered evidence screen had
DeepSeek-V4.1-Flash judge 16{,}032 triples against their CC BY or CC0
source sentence: 56.0\% are supported, 22.6\% partially, 18.2\% not
supported and 3.3\% conflict with the sentence. An earlier evidence
check over the eligible CC BY/CC0 triples (16{,}132 judged) found the
source sentence to state a relationship in 86.6\% of cases but to
support the original label in only 22.0\%; after cleaning-based
relabelling, support rises from 17.2\% to 50.4\%, and kept triples are
more often supported than dropped ones (30.6\% versus 17.2\%).
Against curated references, 13.2\% of 622 resolved microbe--disease
edges appear in Disbiome~\cite{janssens2018disbiome} or BugSigDB
against 5.8\% chance, and kept edges are enriched over dropped ones
(15.5\% versus 6.4\%). BugSigDB support for taxon triples is likewise
higher for kept than dropped triples (54.3\% versus 44.3\%,
Figure~\ref{fig:results}d), a trend that does not survive Holm
correction. Chemical--disease edges resolve less convincingly (11.4\%
in CTD~\cite{davis2026ctd} versus 9.8\% chance). The resulting graph
has 22{,}579 nodes and 25{,}788 edges (cleaned: 8{,}897 and 9{,}682),
with the largest connected component covering 68.5\% (full) and 67.8\%
(cleaned) of nodes.

\section{Discussion}\label{sec:discussion}
The results show that a literature corpus can be turned into a
provenance-traceable common data model whose layers can be improved
without rewriting the frozen builds: re-typing removed the generic entity
type from every kept term and raised the share of triples with a
specific, cue-verified relation to 54.2\% (58.0\% with the triples derived
from catalogue findings), while flagging the triples
it cannot verify instead of forcing a label. The
evaluation also shows where the approach stops: reliability differs by
entity type, model judges agree only moderately, and findings reported
only in figures and supplements are not captured. The recommendations
and limitations below follow from these results.

\subsection{Relation to existing resources}\label{sec:compare}
Microbiome-CDM is designed to sit alongside existing microbiome resources
rather than to replace them (Table~\ref{tab:compare}). Curated databases and
annotated corpora offer expert-level precision and ground truth;
knowledge graphs offer integration, scoring and query; sample-level databases
offer reanalysis of the sequencing data. Microbiome-CDM contributes the
study context and stated findings of the full-text literature, at the level
of the individual sentence and the individual extraction program. The pairings
are concrete. Our coverage of BugSigDB signatures is lowest for those recorded
in supplements (22.1\%), which is exactly the source MicrobiomeKG mines, so
the two are complementary. GutBiomeDB curates the 168,464 samples of the same
compendium, and our model covers the 453 papers behind them, so a sample-level
query can be joined to the sentence that reports the study design. Annotated corpora such as DiMB-RE and MiDRED served as benchmarks for our
relation typing, and our layer in turn offers them a much larger pool of
candidate sentences to annotate. Measured against these resources (Results),
the CDM's study metadata agree closely with curated values, its genus-level
directions agree with two other sources in about 72\% of comparable
cells, and its zero-shot relation typing is precise when it commits to a sign
but below fine-tuned extractors in overall F1.
Because every row names its program, any of these resources can take only the
layers it trusts, for example mappings accepted by two-model agreement or
triples with a verified cue phrase. Its weakest point, compared with curated
resources, is human-annotated ground truth: our human checks (500 mappings,
100 triples) are small samples, not an annotated corpus; we report that gap in
the Limitations.

\subsection{Recommendations}\label{sec:recommendations}
For a new corpus of the same kind, we recommend:

\begin{itemize}
\item Run trained taggers AND zero-shot recognisers side by side, and
choose per entity type, not per system. On BioRED the best organism,
disease and gene recognisers were not trained on BioRED at all, so
performance on the target corpus, not on the training-domain benchmark, should decide.
\item Do not read public benchmark F1 as corpus accuracy: its
correlation with corpus grounding is 0.749 for diseases but 0.023
for organisms. Re-check chosen systems against a sample of the
target corpus before committing to a configuration.
\item Keep each program's mentions separate in the CDM instead of
merging them: when the LLM sources are removed, the majority of
accepted mappings and of terms survive, so provenance per program
supports later re-weighting without re-extraction.
\item Prefer explicit non-values over silent imputation. NOT\_REPORTED
and CONFLICTING labels preserve ambiguity (46{,}182 and 4{,}586 values
in our catalogue) and make downstream filtering a user decision.
\item Screen candidate relation triples against their own source
sentence with an LLM judge before building a graph; cleaning raised
judge-perceived support from 17.2\% to 50.4\% and enriched curated
reference coverage.
\item Attach a character-offset quote to every extracted value from
the start; deterministic re-matching against the anchored quote is
the cheapest check available, and 88.2\% of our quotes are verbatim.
\item Licence-check every text before sending it to an external
model API and record the check; this discipline is what later allows
a licence-filtered evaluation at all.
\end{itemize}

\paragraph{Directions for the next build.} Four measured limitations point
to concrete extensions. First, findings reported only in tables, figures and
supplements are covered far less often than findings in the text (51.4\%,
27.6\% and 22.1\% of table-, figure- and supplement-sourced BugSigDB signatures,
against all text-sourced ones), which is why supplement-based resources such as
MicrobiomeKG complement this one; a table and figure-caption pass is the first
candidate. Second, the closed relation vocabulary misses relations worded
outside it: in the two-annotator check most flagged triples stated a real
relation, so a vocabulary expanded from the corpus's flagged triples, not
from the benchmark corpora, is the second. Third, the catalogue fields for DNA
extraction method and sequencing platform are often rejected by the judge
majority (68.9\% and 56.3\%); both draw on finite vocabularies of kits and
instruments, so a gazetteer pass combined with a field-targeted extractor is the
third. Fourth, on DiMB-RE the relation step typed many entity pairs that the
annotators left unrelated; requiring a syntactic connection between the two
entities before typing would address this, but because the problem was
observed on that benchmark, the rule would need to be fixed in advance and
evaluated on new data. None of these changes the central claim that the
resource is auditable and abstains where it cannot verify a relation.

\subsection{Limitations}\label{sec:limitations}
\begin{itemize}
\item \textbf{Limited human review.} Human review covers one blinded
reviewer's check of 500 randomly sampled accepted mappings (497 rated; second-reviewer
overlap pending) and a blinded two-annotator check of 100 randomly sampled
triples (adjudication pending). These are small samples with wide intervals,
not a human-annotated corpus; all other evaluation is computational (concordance
against chance baselines, judged validity screens and curated-reference
comparisons) and is reported as concordance, never as
accuracy~\cite{chen2025llmbench}.
\item \textbf{Relation quality.} After repair and re-typing,
13{,}969 of the 25{,}788 triples (54.2\%) carry a specific relation with
a verified cue phrase; 8{,}068 are flagged as unverified, and the
two-annotator check shows that many of these state a real relation the
relation vocabulary does not cover, so the flag marks unverified triples, not
extraction errors; 1{,}210
state an association whose direction the paper does not give; a further
1{,}604 carry a specific label whose cue phrase was not found and are not
counted as verified. Re-typed relations, directions and term types are model
output, checked on small constructed test sets (17 of 17 sentences for the
CDM triples and 18 of 20 cases for the catalogue-derived triples) but the 100-triple human check judged each triple as originally extracted,
not its re-typed label, and the 195 triples
whose direction was flagged as reversed have not been reviewed. The 4{,}966
catalogue-derived triples are built by string rules from anchored catalogue
findings and then re-typed with their sentence context; they have no
separate human check. Unsigned associations were resolved by
DeepSeek-V4.1-Flash for CC BY and CC0 papers (243 of 959 resolved) and by a
local Qwen3.5-9B model for all others (162 of 749), so resolution quality may differ between licence groups. Applied zero-shot,
our relation typing is below fine-tuned extractors on DiMB-RE (typed F1 0.25
versus 0.445) and MiDRED (micro-F1 0.57 versus 0.91).
\item \textbf{Judges, not experts.} The three catalogue judges reach
moderate agreement (Cohen's $\kappa$ 0.566 for the best pair, Fleiss'
$\kappa$ 0.471), the record-level summary labels they rate are almost
always judged fair and so carry little information, and the triple screens (two judge models, two rubrics) are not
sanity-checked to the same depth for every judge; concordance with
model judges is not concordance with experts. The judges' majority rejects
68.9\% of stated DNA extraction methods, 56.3\% of sequencing platforms,
40.4\% of body sites and 40.2\% of intervention agents, so these catalogue
fields should be treated as low confidence. Each model step was run once
at temperature 0, and run-to-run variance was not measured.
\item \textbf{Model-version drift.} All model judgements reported here
are tied to the exact model checkpoints used (as of 2026-11): judge
behaviour can change across model updates, and any reproduction must
record the exact model versions, as the Methods do for each screen.
\item \textbf{Documented schema gaps.} In build 3.0.3, definitions found in
the papers cover 2.09\% of terms (3{,}397 of 162{,}409, in 6{,}773 definition
records); the definitions added in build 3.1.0 are model-written; the 53{,}780 kept terms re-typed from the
generic ENTITY type carry model-assigned types and mostly ground to no
vocabulary, and four of their types are proposed extensions; 9 of
14 CDM types have no reference vocabulary; the build declares no
foreign keys (multi-valued links are modelled with explicit link
tables, so referential integrity was checked by audit, not by
declared constraints); 79 change-log entries 3.0.2 $\to$ 3.0.3 were
restored, not recomputed.
\item \textbf{Transfer and scope.} Public-benchmark to corpus
transfer was established for diseases only; for organisms, an
advantage for trained taggers on abstracts was not established
outside their training data. Evaluation is restricted to the 453
gut-microbiome papers linked to one compendium, over 4{,}006
sections; generalisation to other bodies of literature is not
established. The BugSigDB comparison shares only 47 papers with
BugSigDB, and the genus-by-condition map (25 of 44 comparable cells
agree with independent BugSigDB consensus) summarises reported findings;
it is not a meta-analysis and does not establish clinical validity. The
text of the 179 papers whose licences do not permit redistribution cannot
be deposited, so the text-processing steps can be rerun on those papers
only from licensed copies.
\end{itemize}

\section{Conclusions}\label{sec:conclusions}
A multi-program extraction pipeline can turn the literature underlying a
large microbiome compendium into a provenance-traceable common data model
in which every value carries its paper, section, character offsets and
the program that produced it. Curated layers added on top of the frozen
build give every kept term a specific type and 18,935 triples a specific,
cue-verified relation, and three local model judges screen catalogue
values with moderate agreement. Reliability remains uneven across entity
types, and all judge-based results are concordance measures rather than
accuracy against human annotation. The builds are frozen with checksummed
manifests for public deposit.

\backmatter

\section*{Abbreviations}
\begin{list}{}{%
  \setlength{\labelwidth}{6.5em}%
  \setlength{\labelsep}{0.5em}%
  \setlength{\leftmargin}{7em}%
  \setlength{\itemsep}{0pt}%
  \setlength{\parsep}{0pt}%
  \setlength{\topsep}{2pt}%
  \renewcommand{\makelabel}[1]{\textbf{#1}\hfill}}
\item[API] application programming interface
\item[BC5CDR] BioCreative V chemical--disease relation corpus
\item[CDM] common data model
\item[CTD] Comparative Toxicogenomics Database
\item[DOI] digital object identifier
\item[F1] harmonic mean of precision and recall
\item[IBD] inflammatory bowel disease
\item[LLM] large language model
\item[MeSH] Medical Subject Headings
\item[NCBI] National Center for Biotechnology Information
\item[NER] named entity recognition
\item[OMOP] Observational Medical Outcomes Partnership
\item[UMLS] Unified Medical Language System
\end{list}

\section*{Declarations}

\bmhead{Ethics approval and consent to participate}
Not applicable. This study analysed only previously published literature and
public data; no human participants, human tissue or animal subjects were
involved.

\bmhead{Consent for publication}
Not applicable.

\bmhead{Availability of data and materials}
The outputs of this pipeline will be deposited on Zenodo before
publication and made available to reviewers at submission: CDM build
3.1.0 (the unchanged 3.0.3 relational database, the curated 3.1.0 layers
and checksummed manifests), the data catalogue release 3.1.1, the
terminology release 1.0.1, the frozen claims register used for number
verification, and the benchmark and evaluation artefacts, DOI [to be
added when the deposit is published]. The pipeline code will be released
in a public repository with an archived version, [repository and DOI to
be added by the authors]. The underlying source texts are third-party
publications and are redistributed only under their own licences: all
453 papers were licence-checked, and the text of the 179 papers whose
licence does not permit redistribution (135 without a licence and 44
under non-commercial or no-derivatives terms) is withheld from the
deposit and is available from the original publishers. Derived annotations
over those papers (terms, mappings, triples, character offsets and text
hashes) contain no running text, but verbatim evidence quotations do; quotations
are included only for the 274 CC BY or CC0 papers, and the release of derived
annotations for the 135 papers without a licence is under legal review.

\bmhead{Competing interests}
\textbf{[AUTHOR INPUT REQUIRED --- per-author competing-interests statement
missing.} Replace this paragraph with ``The authors declare that they have
no competing interests.'' if that is true for all three authors, or with the
individual disclosures.\textbf{]}

\bmhead{Funding}
\textbf{[AUTHOR INPUT REQUIRED --- funding statement missing.} Replace this
paragraph with the funder and grant numbers, or with ``This research received
no specific grant from any funding agency in the public, commercial or
not-for-profit sectors.''\textbf{]}

\bmhead{Authors' contributions}
\textbf{[AUTHOR INPUT REQUIRED --- CRediT statement missing.} Replace this
paragraph with one CRediT role list per author, drawn from the taxonomy at
\nolinkurl{https://credit.niso.org/}.\textbf{]}

\bmhead{Acknowledgements}

\textbf{[AUTHOR INPUT REQUIRED --- acknowledgements missing.} Replace this
paragraph with the individuals, facilities and compute or storage allocations
to be acknowledged, or with ``Not applicable.'' if there are none.\textbf{]}

\bmhead{Declaration of generative AI and AI-assisted technologies in the
writing and research process}
During the preparation of this work the
authors used generative AI in two distinct roles.
\textbf{(1) As pipeline components and evaluation instruments} (Methods,
Stages 2--5): the small language models Qwen3.5-4B (two extraction runs),
Gemma, Llama and Qwen with few-shot prompting served as entity- and
relation-extraction systems under evaluation; model judges were used as
computational screens, namely DeepSeek-V4-Flash for the string-validity
screen (730{,}195 judged items), DeepSeek-V4.1-Flash, accessed through the
Fraunhofer Genie gateway, for the licence-filtered triple evidence screen
(16{,}032 triples judged against their Creative Commons CC BY or CC0 source
sentences, drawn from 25{,}714 licence-screened sentences), Gemma and
DeepSeek-R1 as the two adjudicating models for vocabulary mapping, and
Qwen3.5-9B, Llama-3.1-8B and Mistral-7B, run locally, as judges of
catalogue values. Where the
manuscript reports a percentage from these screens, the responsible judge
model is named alongside it. All judge-derived results are reported as
concordance or judged validity, not as accuracy against human annotation;
human evaluation consists of the blinded review of 500 sampled accepted
mappings and the blinded two-annotator check of 100 sampled triples. Qwen3.5-9B, run locally, re-typed generic terms and unspecific
triples and typed triples derived from catalogue findings (Methods); to
resolve unsigned associations, DeepSeek-V4.1-Flash on the Genie service
read text from CC~BY or CC0 papers only and the local Qwen3.5-9B read all
other text. For build 3.1.0, Qwen3.5-4B also generated contextual
definitions and repaired triple relations, SapBERT and MiniLM embeddings
supported deduplication, and DeepSeek-V4.1-Flash screened term and concept
labels for relevance.
\textbf{(2) In manuscript and code preparation:} an AI coding assistant
(Claude Code, and opencode with the GLM-5.3 model) was used to develop
pipeline and analysis code, to verify every number in this manuscript
against a frozen claims register, to convert the authors' Methods draft to
this \LaTeX\ template mechanically, and to produce internal drafting
scaffolds that the authors re-wrote in their own words before inclusion;
figures were produced with the same assistance. The authors reviewed and
edited all AI-assisted outputs and take full responsibility for the content
of this published article.

\bibliography{reference}

\end{document}
